\section{Introduction}
\label{sec:introduction}

The Birch and Swinnerton-Dyer conjecture relates the leading term of the
\(L\)-function of an elliptic curve to its rational points, local component
groups, and Tate--Shafarevich group.  Its rank assertion and its exact
leading-term assertion are distinct: a proof that a Heegner point is
nontorsion does not determine the index of that point or the order of the
Tate--Shafarevich group.  This distinction is particularly important at
the prime two, where real components and integral eigenspace indices
contribute to the formula.

We establish the two-primary formula for every elliptic curve over
\(\mathbb Q\) whose usual two-power Selmer group has corank zero or one.
This resolves the two-primary leading-term assertion of the Birch and
Swinnerton-Dyer conjecture in that Selmer-corank range.  It also gives the
pointwise converse from the Selmer-corank bound to analytic rank at most
one.  The theorem concerns the two-primary part of the leading term;
it makes no assertion about higher Selmer corank or the remaining prime
parts.

\subsection{Statement and normalizations}

Let \(E/\mathbb Q\) be an elliptic curve.  We use the Selmer group with
the usual local Kummer conditions, including the real place, and write
\[
 s_2(E)=\corank_{\mathbb Z_2}\Sel_{2^\infty}(E/\mathbb Q).
\]
Let \(\omega_E\) be a global minimal N\'eron differential, let
\(c_\ell(E)\) be the finite Tamagawa numbers, and put
\[
 \Omega_E=\int_{E(\mathbb R)}|\omega_E|,
 \qquad T_E=E(\mathbb Q)_{\tors}.
\]
Thus \(\Omega_E\) integrates over every real component.  For reduced
\(x(P)=a/b\), set \(h_x(P)=\log\max(|a|,|b|)\), with
\(h_x(O)=0\), and define
\[
 H(P)=\lim_{m\to\infty}4^{-m}h_x([2^m]P),\qquad
 B(P,Q)=\frac{H(P+Q)-H(P)-H(Q)}2.
\]
In particular \(B(P,P)=H(P)\).  Here \(H\) is the canonical height
attached to \(2[O]\); the canonical height attached to \([O]\) is
\(H/2\).  The regulator \(\Reg_E\) is the
determinant of \(B\) on a \(\mathbb Z\)-basis of
\(E(\mathbb Q)/T_E\), and is one in rank zero.  Replacing this lattice
by a subgroup of index \(m\) multiplies the regulator by \(m^2\).
This is the regulator convention in
\cite[Section~1]{BuhlerGrossZagier1985}; its compatibility with the
relative Gross--Zagier height is explained in
Section~\ref{sec:classical}. All constant valuations are normalized by
\(v_2(2)=1\).

\begin{theorem}\label{thm:main}
Let \(E/\mathbb Q\) be any elliptic curve with \(s_2(E)\leq1\).
Then
\[
 r:=\operatorname{rank}E(\mathbb Q)
   =\operatorname{ord}_{s=1}L(E,s)=s_2(E),
 \qquad \#\Sha(E/\mathbb Q)<\infty.
\]
Moreover,
\[
 Q_E=\frac{L^{(r)}(E,1)\,(\#T_E)^2}
 {r!\,\Omega_E\Reg_E\prod_{\ell\ \mathrm{finite}}c_\ell(E)}
 \in\mathbb Q_{>0},
 \qquad
 v_2(Q_E)=v_2\bigl(\#\Sha(E/\mathbb Q)\bigr).
\]
There is no restriction on \(E[2]\) or on the reduction of \(E\) at two.
\end{theorem}

\begin{corollary}[An any-prime input criterion]\label{cor:any-prime-bsd}
Let \(E/\mathbb Q\) be an elliptic curve. Suppose that, for some prime
\(p\), its full \(p\)-power Selmer group with the usual local Kummer
conditions at every place satisfies
\[
 s_p(E):=\corank_{\mathbb Z_p}\Sel_{p^\infty}(E/\mathbb Q)\le1.
\]
Then all conclusions of Theorem~\ref{thm:main} hold, with common rank
\(r=s_p(E)\).
\end{corollary}

\begin{proof}
By \cite[Theorem~1.1]{OpenAIAllPrimeSelmer2026},
\(\operatorname{rank}E(\mathbb Q)=r:=s_p(E)\) and the whole group
\(\Sha(E/\mathbb Q)\) is finite. The Kummer exact sequence at two is
\[
 0\longrightarrow E(\mathbb Q)\otimes(\mathbb Q_2/\mathbb Z_2)
 \longrightarrow\Sel_{2^\infty}(E/\mathbb Q)
 \longrightarrow\Sha(E/\mathbb Q)[2^\infty]\longrightarrow0.
\]
Taking coranks gives \(s_2(E)=\operatorname{rank}E(\mathbb Q)=r\le1\),
so Theorem~\ref{thm:main} applies.
\end{proof}

There is also a density consequence of Theorem~\ref{thm:main}.
Combining it with the signed-squarefree Selmer-corank densities in
\cite[proof of Theorem~1.2]{OpenAIGoldfeldDensity2026} gives the exact
two-primary formula for a density-one set of quadratic twists of each
fixed \(E/\mathbb Q\). The common analytic and Mordell--Weil rank is
zero or one, with each value having density \(1/2\); see
Corollary~\ref{cor:quadratic-twist-density}.

\subsection{History and the integral difficulty}

The conjecture grew out of the computations of Birch and
Swinnerton-Dyer~\cite{BSD65}. The modularity theorem developed from
the semistable case of Wiles and Taylor--Wiles~\cite{Wiles95,TaylorWiles95}
to all elliptic curves over \(\mathbb Q\)~\cite{BCDT01}.
Together with the work of Gross--Zagier and Kolyvagin, it shows that
analytic rank zero or one implies the
corresponding rational rank and finiteness of the whole
Tate--Shafarevich group~\cite{BCDT01,GZ86,Kolyvagin90}. These analytic
low-rank theorems are inputs here. The converse from the full two-power
Selmer group, and the exact leading coefficient at two, require
additional integral information. Cassels established the isogeny
comparison for elliptic curves; Tate formulated its abelian-variety
extension, including the restriction-of-scalars comparison used
below~\cite{Cassels65,Tate1966}.

At odd primes, exact leading-term formulas in analytic ranks zero and
one were established under residual and local hypotheses by
Skinner--Urban, Wei Zhang, and Jetchev--Skinner--Wan, through Iwasawa
main conjectures and the integral theory of Heegner points
\cite{SkinnerUrban2014,Zhang2014,JetchevSkinnerWan2017}.

The CM theory provides a substantial part of the exact-formula
history. Elliptic units and Rubin's imaginary-quadratic main conjectures
supply its Iwasawa-theoretic background~\cite{Rubin1991}.
Johnson-Leung--Kings establish the individual-character main conjecture
without prime exceptions~\cite{JLK11}; Burungale--Flach use this input
to prove full BSD for every CM elliptic curve over \(\mathbb Q\) of
analytic rank zero~\cite{BF23}. Burungale--Tian's rank-zero
\(p\)-converse, valid for every prime, then shows that the CM
Selmer-corank-zero case already follows from these results
\cite{BurungaleTian26}.

For CM curves of analytic rank one, Li--Tian--Yan--Zhu prove the exact
two-primary formula under good ordinary reduction at two. Their
Theorem~1.2 combines this formula with earlier ordinary CM
two-converse results, including the work of Burungale--Castella--Skinner--Tian
\cite{LTYZ2025,BCST22}. The direct CM comparison below also treats
Selmer corank one at every other reduction type at two, and supplies
the paired determinant comparison used in the non-CM argument.

For non-CM curves, and for explicit CM twist families, the prime two
has required separate methods. Zhao developed lower bounds and
combinatorial equality criteria for two-adic central-value
valuations~\cite{Zhao1997}. Coates--Li--Tian--Zhai prove full BSD for
specified rank-zero twists of \(X_0(49)\), together with rank-one
and finiteness results for other restricted twists~\cite{CLTZ2015}.
Cai--Li--Zhai propagate the exact two-part in rank-zero twist
families under base-curve and local hypotheses~\cite{CLZ2020}.
Shu--Zhai obtain exact rank-zero and rank-one two-part formulas for
families satisfying rational two-isogeny, base-formula, and splitting
conditions~\cite{ShuZhai21}. These results establish important
special cases. For curves without rational two-torsion, Kriz--Li
propagate the exact two-part in rank-zero and rank-one twist families
under a Heegner-logarithm unit condition and further local and
base-formula hypotheses~\cite[Theorem~5.1]{KrizLi2019}.
The present theorem permits every residual image and every reduction
type at two.

Kato's zeta elements and explicit reciprocity law connect integral
Galois cohomology to modular periods~\cite{Kato04}. The zeta-element
constructions of Burungale--Skinner--Tian--Wan give further low-rank
applications with prime and local hypotheses~\cite{BSTW24}.
Our proof uses Kato's full-level reciprocity law in a precise
weight-two specialization, while proving the integral primitivity
statements required at two. The logarithm-squared comparison belongs
to the framework proposed by Perrin-Riou and developed at odd
semistable primes by Bertolini--Darmon--Venerucci
\cite{PerrinRiou1993,BertoliniDarmonVenerucci2022}. Depleted CM
primitives and square-logarithm formulas have their antecedents in
Bertolini--Darmon--Prasanna~\cite{BertoliniDarmonPrasanna2013}.
Here these comparisons are combined with separate integral arguments
at two.

Determinant lines and the Tamagawa-number
formulation provide the language for these comparisons
\cite{KnudsenMumford76,BlochKato90}; Selmer complexes and Bockstein
pairings organize their first-order variation~\cite{Nekovar06}.
Finite--singular comparisons also form a basic part of the theory of
Kolyvagin systems~\cite{MazurRubin04}. The finite models, exact
switch volumes, and dyadic alternation needed here are proved below,
with explicit treatment of real cohomology.

The integral difficulty is easy to describe.  An Euler-system
construction may produce a point or a determinant generator only after
multiplication by a fixed nonzero integer.  Such a bound can prove rank
and finiteness, but its unknown two-adic valuation cannot determine
\(v_2(Q_E)\).  We therefore separate two operations.  Horizontal
comparisons, made at height-one primes where two is invertible, allow
fixed clearing denominators. Residual concentration supplies the
remaining integrality, and an independent primitivity argument is
needed at the closed point. For example, for a rational coordinate
\(U\) over \(\mathbb Z_2[[u]]\), nonnegative orders at every
height-one prime imply \(U\in\mathbb Z_2[[u]]\). The stronger
conclusion that \(U\) is a unit requires its constant term to be odd.
Our determinant comparisons keep these steps separate and retain the
local Euler, real-component, and lattice factors through specialization.

\subsection{The three comparisons}

Write \(\an(A)=\operatorname{ord}_{s=1}L(A,s)\), and, when
\(\an(A)\leq1\), define
\[
 X(A)=v_2\!\left(
 \frac{L^{(\an(A))}(A,1)\,(\#A(\mathbb Q)_{\tors})^2}
 {\an(A)!\,\Omega_A\Reg_A\prod_qc_q(A)\,\#\Sha(A/\mathbb Q)}
 \right).
\]
The classical analytic low-rank theorem makes \(X(A)\) well-defined.
For a quadratic twist parameter \(a\), write \(A^a\), allowing
\(a=1\), and write \(N_A\) for the conductor.  The proof is organized
around the following assertions.

\begin{proposition}[Positive comparison]\label{prop:positive}
Let \(A/\mathbb Q\) be non-CM.  If \(\an(A)\leq1\), then
\(X(A)\geq0\).  Suppose a positive fundamental discriminant \(a\)
prime to \(2N_A\), or \(a=1\), satisfies
\(\an(A^a)\leq1\) and \(X(A^a)=0\).  Then \(s_2(A)\leq1\)
implies \(\an(A)=s_2(A)\) and \(X(A)=0\).
\end{proposition}

\begin{proposition}[Split-pair anchor]\label{prop:anchor}
Let \(A/\mathbb Q\) be non-CM.  There exist a fundamental
discriminant \(h\) prime to \(2N_A\), allowing \(h=1\), and an
imaginary quadratic field of discriminant \(k\) prime to \(2hN_A\),
in which all primes dividing \(2hN_A\) split, such that
\[
 \an(A^h)+\an(A^{hk})=1,
 \qquad X(A^h)+X(A^{hk})=0.
\]
\end{proposition}

\begin{proposition}[CM comparison]\label{prop:cm}
If \(A/\mathbb Q\) has complex multiplication and \(s_2(A)\leq1\),
then \(\an(A)=s_2(A)\) and \(X(A)=0\).
\end{proposition}

These propositions have complementary roles. The anchor supplies two
twists with a vanishing sum of discrepancies; nonnegativity makes each
discrepancy zero. Since \(k<0\), one member of the pair is an allowed
positive twist, providing the input for the transfer in
Proposition~\ref{prop:positive}. The CM comparison treats the remaining
curves directly. Section~\ref{sec:conclusion} completes this deduction,
including the classical rank, finiteness, and rationality inputs.

\subsection{Proof structure and reusable constructions}

The positive comparison uses the actual integral relative-symbol
lattice of a modular elliptic quotient, an unnormalized Siegel-unit
period law, and determinant-preserving finite/singular line switches.
A tame height calculation gives the rank-one logarithmic identity
without an ordinary-reduction hypothesis.  At an initially unknown
center, a Selmer Bockstein detects the first derivative before analytic
rank is assumed.

For the split-pair anchor, finite two-descent supplies a twist and a
split imaginary quadratic field over which the usual Selmer corank is
one. We compare a determinant strict at one place above two with a
product of two depleted CM-disk logarithmic measures. Ring-class
derivatives provide square switches. An integral unit comparison,
together with the independently obtained corank-one condition, makes
the central Heegner trace nontorsion and yields the exact split formula.

The residual primitivity argument depends on the residual image: scalar
constituents for rational two-torsion, a CM comparator for imaginary
\(S_3\), ray units for a cyclic cubic image, and a real-component
detector for real \(S_3\).  The last argument uses an alternating
determinant and a Pfaffian cofactor lift divisible by a product of local
factors, even when those factors are not coprime.

Several constructions apply beyond their immediate branches: marked
finite free cochain models that retain specialization and cup products;
exact line and square switches with specified determinant volumes;
bounded-series division from high-character tests; conjugate Bockstein
alternation at two; and the integral Pfaffian cofactor lift.  We give
the hypotheses of these constructions explicitly.  Figure~\ref{fig:logic}
records the part of the proof order that prevents the unit and rank
arguments from being used circularly.

\begin{figure}[htbp]
\centering
\begin{tikzpicture}[
 box/.style={draw,rounded corners=2pt,align=center,font=\small,
 text width=4.9cm,inner sep=6pt},
 >={Stealth[length=2mm]},node distance=10mm and 10mm]
 \node[box] (base) at (0,0) {Finite models, exact periods,\\horizontal derivatives};
 \node[box] (positive) at (-3.25,-1.8) {Positive comparison:\\nonnegativity and unit transfer};
 \node[box] (split) at (3.25,-1.8) {Split comparison:\\Selmer corank one and a unit\\give a nonzero central point};
 \node[box] (anchor) at (3.25,-3.8) {Finite Selmer seeds and\\residual primitivity:\\split-pair anchors};
 \node[box] (target) at (-3.25,-4.5) {A positive unit anchor;\\transfer to the original curve};
 \draw[->] (base) -- (positive);
 \draw[->] (base) -- (split);
 \draw[->] (split) -- (anchor);
 \draw[->] (positive) -- (target);
 \draw[->] (anchor) -- (target);
\end{tikzpicture}
\caption{Logical order in the non-CM proof.  Nonnegativity precedes
the anchor.  The Hecke-field rank test used by the comparators starts
with a classically known simple analytic product and does not use a
residual unit conclusion.  The direct CM argument is treated separately.}
\label{fig:logic}
\end{figure}

\paragraph{Navigation.}
Section~\ref{sec:classical} fixes the classical inputs and exact
Gross--Zagier factors; Section~\ref{sec:foundations} proves the common
cochain and determinant lemmas. Sections~\ref{sec:siegel}--\ref{sec:positive-completion}
prove Proposition~\ref{prop:positive}, with the unknown-rank
Bockstein argument in Section~\ref{subsec:positive-completion:bockstein}.
Sections~\ref{sec:split-measures}--\ref{sec:split-local} establish the
split determinant, its horizontal divisibility, and the central rank
test. The reducible anchor is completed in Section~\ref{sec:reducible}.
Section~\ref{sec:cm} proves Proposition~\ref{prop:cm} and the CM
comparison used later. After the alternation calculation in
Section~\ref{sec:alternation-imaginary}, the imaginary \(S_3\), real
\(S_3\), and cyclic-cubic anchors are proved in
Sections~\ref{alternation-imaginary:anchor-section}, \ref{sec:real},
and~\ref{sec:cyclic}, respectively. Section~\ref{sec:conclusion}
assembles these branches.

Our notation is local to each construction when a field, coefficient
ring, or modular factor changes.  Orders at height-one primes refer to
the corresponding discrete valuation rings; valuations of dyadic
constants always retain the normalization \(v_2(2)=1\).
